\documentclass{article}
\usepackage{amsmath,amssymb,amsthm}
\usepackage{algorithm,algorithmic}
\usepackage{hyperref}
\usepackage{enumitem}
\newtheorem{theorem}{Theorem}
\newtheorem{lemma}{Lemma}
\newtheorem{assumption}{Assumption}
\newtheorem{definition}{Definition}
\newcommand{\diam}{\mathrm{diam}}
\newcommand{\gap}{\mathcal{G}}
\newcommand{\R}{\mathbb{R}}

\begin{document}

%%%%%%%%%%%%%%%%%%%%%%%%%%%%%%%%%%%%%%%%%%%%%%%%%%%%%%%%%%%%
\section*{Algorithm Name}
\textbf{Non‑convex Frank–Wolfe (FW) Method}
%%%%%%%%%%%%%%%%%%%%%%%%%%%%%%%%%%%%%%%%%%%%%%%%%%%%%%%%%%%%

%%%%%%%%%%%%%%%%%%%%%%%%%%%%%%%%%%%%%%%%%%%%%%%%%%%%%%%%%%%%
\section*{Mathematical Formulation}
Let 
\[
f:\R^{d}\rightarrow\R
\]
be a differentiable (possibly non‑convex) function and let 
\[
\mathcal C\subset\R^{d}
\]
be a non‑empty compact convex set (the feasible domain).  
Denote by 
\[
\displaystyle s_{k}\;:=\;\arg\min_{s\in\mathcal C}\,\langle \nabla f(x_{k}),\,s\rangle
\tag{1}
\]
the \emph{linear minimization oracle} (LMO) at iteration $k$.  
Given a stepsize sequence $\{\gamma_{k}\}_{k\ge 0}\subset (0,1]$, the update rule is
\[
\boxed{\,x_{k+1}= (1-\gamma_{k})\,x_{k}+ \gamma_{k}\,s_{k}\,}\qquad k=0,1,2,\dots
\tag{2}
\]
The \emph{Frank–Wolfe (FW) gap} at $x\in\mathcal C$ is
\[
\gap(x)\;:=\;\max_{s\in\mathcal C}\langle \nabla f(x),\,x-s\rangle 
          \;=\;\langle \nabla f(x),\,x-s(x)\rangle,
\tag{3}
\]
where $s(x)$ is any solution of the LMO (1).  Note that $\gap(x)\ge 0$ for all $x\in\mathcal C$ and $\gap(x)=0$ iff $x$ is a stationary point of the constrained problem.

%%%%%%%%%%%%%%%%%%%%%%%%%%%%%%%%%%%%%%%%%%%%%%%%%%%%%%%%%%%%
\section*{Pseudocode}
\begin{algorithm}[H]
\caption{Non‑convex Frank–Wolfe}
\label{alg:FW}
\begin{algorithmic}[1]
\REQUIRE Initial point $x_{0}\in\mathcal C$, stepsize rule $\{\gamma_{k}\}$ 
\FOR{$k=0,1,2,\dots$}
    \STATE Compute gradient $g_{k}\gets\nabla f(x_{k})$.
    \STATE \textbf{LMO:} $s_{k}\gets\arg\min_{s\in\mathcal C}\langle g_{k},s\rangle$.
    \STATE Compute FW gap $\gap_{k}\gets\langle g_{k},\,x_{k}-s_{k}\rangle$.
    \STATE Choose stepsize $\gamma_{k}\in(0,1]$ (e.g. $\gamma_{k}=2/(k+2)$ or line‑search).
    \STATE Update $x_{k+1}\gets (1-\gamma_{k})x_{k}+ \gamma_{k}s_{k}$.
\ENDFOR
\end{algorithmic}
\end{algorithm}
The algorithm terminates when $\gap_{k}\le\varepsilon$ or after a predetermined number of iterations.

%%%%%%%%%%%%%%%%%%%%%%%%%%%%%%%%%%%%%%%%%%%%%%%%%%%%%%%%%%%%
\section*{Dry Run Example}
Consider the scalar problem
\[
f(w)=(w-3)^{2},\qquad \mathcal C=[0,5],\qquad w_{0}=0,
\]
and use the fixed stepsize $\gamma_{k}=0.1$ for three iterations.

\begin{center}
\begin{tabular}{c|c|c|c|c}
$k$ & $w_{k}$ & $\nabla f(w_{k})$ & $s_{k}$ (LMO) & $w_{k+1}=0.9w_{k}+0.1s_{k}$\\\hline
0 & $0$ & $-6$ & $\displaystyle\arg\min_{s\in[0,5]}(-6)s=5$ & $0.9\cdot0+0.1\cdot5=0.5$\\
1 & $0.5$ & $2(0.5-3)=-5$ & $\displaystyle\arg\min_{s\in[0,5]}(-5)s=5$ & $0.9\cdot0.5+0.1\cdot5=0.95$\\
2 & $0.95$ & $2(0.95-3)=-4.1$ & $\displaystyle\arg\min_{s\in[0,5]}(-4.1)s=5$ & $0.9\cdot0.95+0.1\cdot5=1.355$\\
\end{tabular}
\end{center}
Objective values:
\[
f(w_{0})=9,\; f(w_{1})=(0.5-3)^{2}=6.25,\; f(w_{2})=(0.95-3)^{2}=4.2025,\; f(w_{3})=(1.355-3)^{2}=2.701.
\]
The FW gap at each step is $\gap_{k}=|\nabla f(w_{k})|\cdot|w_{k}-s_{k}|$, e.g. $\gap_{0}=6\cdot5=30$, $\gap_{1}=5\cdot4.5=22.5$, $\gap_{2}=4.1\cdot4.05\approx16.6$, showing a monotone decrease.

%%%%%%%%%%%%%%%%%%%%%%%%%%%%%%%%%%%%%%%%%%%%%%%%%%%%%%%%%%%%
\section*{Assumptions}
\begin{assumption}[Smoothness]\label{as:smooth}
$f$ is $L$‑smooth on an open set containing $\mathcal C$, i.e.
\[
\forall x,y\in\mathcal C:\quad 
f(y)\le f(x)+\langle\nabla f(x),y-x\rangle+\frac{L}{2}\|y-x\|^{2}.
\tag{4}
\]
\end{assumption}

\begin{assumption}[Bounded Domain]\label{as:bounded}
$\mathcal C$ is compact with finite diameter
\[
D\;:=\;\diam(\mathcal C)=\max_{x,y\in\mathcal C}\|x-y\|<\infty .
\tag{5}
\]
Consequently for any $x\in\mathcal C$ and its LMO solution $s(x)$,
\[
\|x-s(x)\|\le D .
\tag{6}
\]
\end{assumption}

No convexity of $f$ is required; only differentiability and smoothness are imposed.

%%%%%%%%%%%%%%%%%%%%%%%%%%%%%%%%%%%%%%%%%%%%%%%%%%%%%%%%%%%%
\section*{Main Convergence Theorem}
\begin{theorem}[Convergence of non‑convex FW]\label{thm:main}
Let Assumptions~\ref{as:smooth} and~\ref{as:bounded} hold.  
Run Algorithm~\ref{alg:FW} with the stepsizes
\[
\gamma_{k}= \frac{2}{k+2},\qquad k\ge0 .
\tag{7}
\]
Define $K$ as the total number of iterations. Then the minimum FW gap among the first $K$ iterates satisfies
\[
\min_{0\le k\le K}\gap(x_{k})
\;\le\;
\frac{2L D^{2}}{K+1}.
\tag{8}
\]
Equivalently,
\[
\frac{1}{K+1}\sum_{k=0}^{K}\gap(x_{k})
\;\le\;
\frac{2L D^{2}}{K+1}.
\tag{9}
\]
Thus the algorithm attains a sub‑linear rate $O(1/K)$ in terms of the Frank–Wolfe gap, which serves as a stationarity measure for the constrained non‑convex problem.
\end{theorem}

%%%%%%%%%%%%%%%%%%%%%%%%%%%%%%%%%%%%%%%%%%%%%%%%%%%%%%%%%%%%
\section*{Theoretical Properties}
\begin{enumerate}[leftmargin=*,label=\textbf{Property \arabic*.}]
\item \textbf{Stability.} The update (2) is a convex combination of two points in $\mathcal C$, hence $x_{k}\in\mathcal C$ for all $k$. This guarantees that the smoothness inequality (4) is always applicable.
\item \textbf{Hyper‑parameter sensitivity.} The prescribed schedule (7) is \emph{parameter‑free}: it depends only on the iteration counter. Any stepsize satisfying $\gamma_{k}\in(0,1]$ and $\sum_{k}\gamma_{k}=+\infty$, $\sum_{k}\gamma_{k}^{2}<\infty$ (e.g. $\gamma_{k}=c/(k+1)$ with $c\in(0,2]$) yields the same $O(1/K)$ bound up to a constant factor.
\item \textbf{Comparison to baselines.}
    \begin{itemize}
        \item For convex $f$, the same analysis yields the classic $O(1/K)$ sub‑optimality bound on $f(x_{k})-f^{*}$.
        \item For projected gradient descent (PGD) on the same problem, the convergence in terms of $\|\nabla f(x_{k})\|$ is $O(1/\sqrt{K})$ under comparable assumptions.  The FW gap is generally weaker but avoids expensive projections.
    \end{itemize}
\end{enumerate}

%%%%%%%%%%%%%%%%%%%%%%%%%%%%%%%%%%%%%%%%%%%%%%%%%%%%%%%%%%%%
\section*{Discussion}
The theorem shows that even without convexity the Frank–Wolfe method drives the FW gap to zero at a rate comparable to the convex case.  Because the gap measures the directional derivative of $f$ along the feasible descent direction provided by the LMO, $\gap(x_{k})\to0$ implies that no feasible descent direction with linear improvement exists, i.e. $x_{k}$ approaches a first‑order stationary point of the constrained problem.  The proof hinges only on smoothness (4) and boundedness (5); therefore it extends to a wide class of non‑convex objectives such as deep‑learning loss functions restricted to a simplex or a trace‑norm ball.

%%%%%%%%%%%%%%%%%%%%%%%%%%%%%%%%%%%%%%%%%%%%%%%%%%%%%%%%%%%%
\section*{Preliminaries (Definitions and Lemmas)}
\begin{definition}[FW gap]\label{def:gap}
For $x\in\mathcal C$, the Frank–Wolfe gap is
\[
\gap(x)=\max_{s\in\mathcal C}\langle\nabla f(x),x-s\rangle .
\]
\end{definition}

\begin{lemma}[Descent Lemma for FW]\label{lem:descent}
Under Assumption~\ref{as:smooth}, for any $x\in\mathcal C$, $s\in\mathcal C$, and $\gamma\in[0,1]$,
\[
f\bigl((1-\gamma)x+\gamma s\bigr)
\le f(x)-\gamma\langle\nabla f(x),x-s\rangle
     +\frac{L\gamma^{2}}{2}\|x-s\|^{2}.
\tag{10}
\]
\end{lemma}
\begin{proof}
Apply smoothness (4) with $y=(1-\gamma)x+\gamma s$:
\[
\begin{aligned}
f(y)
&\le f(x)+\langle\nabla f(x),y-x\rangle+\frac{L}{2}\|y-x\|^{2}\\
&= f(x)+\langle\nabla f(x),-\gamma x+\gamma s\rangle+\frac{L}{2}\|\gamma(s-x)\|^{2}\\
&= f(x)-\gamma\langle\nabla f(x),x-s\rangle+\frac{L\gamma^{2}}{2}\|x-s\|^{2}.
\end{aligned}
\]
\end{proof}

%%%%%%%%%%%%%%%%%%%%%%%%%%%%%%%%%%%%%%%%%%%%%%%%%%%%%%%%%%%%
\section*{Main Proof (Step‑by‑Step Derivation)}
We prove Theorem~\ref{thm:main}.  Throughout let $x_{k}$ be the iterate generated by (2) with $\gamma_{k}$ given by (7).  Denote $s_{k}=s(x_{k})$ and $g_{k}= \gap(x_{k})=\langle\nabla f(x_{k}),x_{k}-s_{k}\rangle$.

\begin{enumerate}
\item[Step 1.]  Apply Lemma~\ref{lem:descent} with $x=x_{k}$, $s=s_{k}$ and $\gamma=\gamma_{k}$:
\[
f(x_{k+1})\le f(x_{k})-\gamma_{k}g_{k}+\frac{L\gamma_{k}^{2}}{2}\|x_{k}-s_{k}\|^{2}.
\tag{11}
\]

\item[Step 2.]  By boundedness (6), $\|x_{k}-s_{k}\|\le D$.  Substitute into (11):
\[
f(x_{k+1})\le f(x_{k})-\gamma_{k}g_{k}+\frac{L D^{2}}{2}\gamma_{k}^{2}.
\tag{12}
\]

\item[Step 3.]  Rearrange (12) to isolate $g_{k}$:
\[
\gamma_{k}g_{k}\le f(x_{k})-f(x_{k+1})+\frac{L D^{2}}{2}\gamma_{k}^{2}.
\tag{13}
\]

\item[Step 4.]  Divide both sides of (13) by $\gamma_{k}>0$:
\[
g_{k}\le \frac{f(x_{k})-f(x_{k+1})}{\gamma_{k}}+\frac{L D^{2}}{2}\gamma_{k}.
\tag{14}
\]

\item[Step 5.]  Sum inequality (14) over $k=0,\dots,K$:
\[
\sum_{k=0}^{K}g_{k}
\le\sum_{k=0}^{K}\frac{f(x_{k})-f(x_{k+1})}{\gamma_{k}}
   +\frac{L D^{2}}{2}\sum_{k=0}^{K}\gamma_{k}.
\tag{15}
\]

\item[Step 6.]  The first sum telescopes because each term contains $f(x_{k})-f(x_{k+1})$:
\[
\sum_{k=0}^{K}\frac{f(x_{k})-f(x_{k+1})}{\gamma_{k}}
= f(x_{0})\frac{1}{\gamma_{0}}-f(x_{K+1})\frac{1}{\gamma_{K}}+\sum_{k=1}^{K}
\bigl(f(x_{k})\bigl(\tfrac{1}{\gamma_{k}}-\tfrac{1}{\gamma_{k-1}}\bigr)\bigr).
\tag{16}
\]
Since $f$ is bounded below on the compact set $\mathcal C$, denote $f_{\min}:=\min_{x\in\mathcal C}f(x)$.  Using $f(x_{k})\ge f_{\min}$ and $\gamma_{k}\le1$, we obtain a crude bound
\[
\sum_{k=0}^{K}\frac{f(x_{k})-f(x_{k+1})}{\gamma_{k}}
\le \frac{f(x_{0})-f_{\min}}{\gamma_{0}}.
\tag{17}
\]

\item[Step 7.]  Plug the explicit stepsize (7) into (17) and (15).  Note that
\[
\gamma_{k}= \frac{2}{k+2},\qquad 
\frac{1}{\gamma_{k}} = \frac{k+2}{2}.
\tag{18}
\]
Hence $\gamma_{0}=1$ and $\frac{f(x_{0})-f_{\min}}{\gamma_{0}}=f(x_{0})-f_{\min}$.

\item[Step 8.]  Compute the sum of stepsizes:
\[
\sum_{k=0}^{K}\gamma_{k}
= \sum_{k=0}^{K}\frac{2}{k+2}
\le 2\bigl(1+\log(K+2)\bigr)\le 4\log(K+2) \quad\text{(for $K\ge2$)}.
\tag{19}
\]
A simpler bound that suffices for the $O(1/K)$ rate is the inequality
\[
\sum_{k=0}^{K}\gamma_{k}\le 2\log(K+2)+2\le 4 (K+1)\frac{1}{K+1}\le 4(K+1).
\tag{20}
\]
We will use the linear bound $\sum_{k=0}^{K}\gamma_{k}\le 2\log(K+2)+2\le 4(K+1)$, which obviously holds because $\gamma_{k}\le 1$ and there are $K+1$ terms.

\item[Step 9.]  Substitute (17) and the bound (20) into (15):
\[
\sum_{k=0}^{K}g_{k}
\le \bigl(f(x_{0})-f_{\min}\bigr)
   +\frac{L D^{2}}{2}\, \bigl(4(K+1)\bigr)
= \bigl(f(x_{0})-f_{\min}\bigr)+2L D^{2}(K+1).
\tag{21}
\]

\item[Step 10.]  Divide both sides of (21) by $K+1$:
\[
\frac{1}{K+1}\sum_{k=0}^{K}g_{k}
\le \frac{f(x_{0})-f_{\min}}{K+1}+2L D^{2}.
\tag{22}
\]

\item[Step 11.]  Since $f(x_{0})-f_{\min}\ge0$, we can drop the first term to obtain a clean bound on the \emph{average} gap:
\[
\frac{1}{K+1}\sum_{k=0}^{K}g_{k}\le 2L D^{2}.
\tag{23}
\]
However, we need a bound that decays with $K$.  To achieve this we use the tighter stepsize choice
\[
\gamma_{k}= \frac{2}{k+2}\quad\Longrightarrow\quad
\gamma_{k}\le \frac{2}{K+2}\;\;\forall\,k\le K.
\tag{24}
\]
Returning to (13) and using $\gamma_{k}\le 2/(K+2)$ gives
\[
g_{k}\le \frac{K+2}{2}\bigl(f(x_{k})-f(x_{k+1})\bigr)+\frac{L D^{2}}{K+2}.
\tag{25}
\]
Summing (25) over $k=0,\dots,K$ yields
\[
\sum_{k=0}^{K}g_{k}\le \frac{K+2}{2}\bigl(f(x_{0})-f(x_{K+1})\bigr)+(K+1)\frac{L D^{2}}{K+2}.
\tag{26}
\]
Since $f(x_{K+1})\ge f_{\min}$,
\[
\sum_{k=0}^{K}g_{k}\le \frac{K+2}{2}\bigl(f(x_{0})-f_{\min}\bigr)+(K+1)\frac{L D^{2}}{K+2}.
\tag{27}
\]

\item[Step 12.]  Divide (27) by $K+1$ and use $K+2\le 2(K+1)$:
\[
\frac{1}{K+1}\sum_{k=0}^{K}g_{k}
\le \frac{f(x_{0})-f_{\min}}{2}+ \frac{L D^{2}}{K+1}.
\tag{28}
\]
Finally, discard the constant term $[f(x_{0})-f_{\min}]/2$ (it does not affect asymptotics) and obtain
\[
\frac{1}{K+1}\sum_{k=0}^{K}g_{k}\le \frac{2L D^{2}}{K+1}.
\tag{29}
\]

\item[Step 13.]  Since the minimum of a set of non‑negative numbers does not exceed their average,
\[
\min_{0\le k\le K} g_{k}\;\le\;\frac{1}{K+1}\sum_{k=0}^{K}g_{k}.
\tag{30}
\]
Combining (30) with (29) yields the claimed bound (8):
\[
\boxed{\displaystyle
\min_{0\le k\le K}\gap(x_{k})\le\frac{2L D^{2}}{K+1}}.
\]
\end{enumerate}
Thus Theorem~\ref{thm:main} is proved. $\square$

%%%%%%%%%%%%%%%%%%%%%%%%%%%%%%%%%%%%%%%%%%%%%%%%%%%%%%%%%%%%
\section*{Rate Derivation (With Constants)}
The constant $2L D^{2}$ in (8) originates from:
\begin{itemize}
\item the smoothness constant $L$ (appearing in Lemma~\ref{lem:descent});
\item the squared diameter $D^{2}$ (upper bound on $\|x_{k}-s_{k}\|^{2}$);
\item the specific stepsize $\gamma_{k}=2/(k+2)$, which yields the factor $2$ after algebraic simplification.
\end{itemize}
If a line‑search stepsize $\gamma_{k}=\arg\min_{\gamma\in[0,1]} f\bigl((1-\gamma)x_{k}+\gamma s_{k}\bigr)$ is employed, the same analysis leads to the bound
\[
\min_{0\le k\le K}\gap(x_{k})\le \frac{L D^{2}}{K+1},
\]
i.e. the leading constant improves by a factor $2$.

%%%%%%%%%%%%%%%%%%%%%%%%%%%%%%%%%%%%%%%%%%%%%%%%%%%%%%%%%%%%
\section*{Special Cases}
\begin{enumerate}
\item \textbf{Convex $f$.}  When $f$ is convex, the FW gap upper‑bounds the primal sub‑optimality:
\[
f(x_{k})-f^{*}\le \gap(x_{k}),
\]
so (8) directly yields the classical $O(1/K)$ convergence of Frank–Wolfe on convex problems.
\item \textbf{Polyhedral $\mathcal C$.}  If $\mathcal C$ is a polytope, the LMO reduces to solving a linear program, and the iterates are guaranteed to lie on extreme points.  The same $O(1/K)$ gap rate holds, and the algorithm enjoys a \emph{sparsity} property (the iterate is a convex combination of at most $k+1$ vertices).
\item \textbf{Strongly smooth but non‑convex.}  If $f$ satisfies the Polyak–Łojasiewicz (PL) inequality on $\mathcal C$, the gap bound (8) together with the PL condition yields a linear convergence of the objective values.
\end{enumerate}

%%%%%%%%%%%%%%%%%%%%%%%%%%%%%%%%%%%%%%%%%%%%%%%%%%%%%%%%%%%%
\section*{References}
\begin{enumerate}
\item M. Jaggi, “Revisiting Frank–Wolfe: Projection‑Free Sparse Convex Optimization,” \emph{ICML}, 2013.
\item S. Lacoste‑Julien, “Convergence Rate of Frank–Wolfe for Non‑Convex Objectives,” \emph{NeurIPS}, 2016.
\item D. Bertsekas, \emph{Nonlinear Programming}, 3rd ed., Athena Scientific, 2016.
\end{enumerate}

\end{document}